\begin{definition}
Let $\Omega\subset\mathbb{C}$ be open, let $k\ge0$ be an integer and let
$0<\alpha\le1$.  A function $u\colon\mathbb{C}\to\mathbb{C}$ belongs to the
H\"older space $C^{k,\alpha}_{\mathrm{loc}}(\Omega)$ if
\begin{enumerate}
\item[(a)] $u$ is of class $C^k$ on $\Omega$, and
\item[(b)] for every compact $F\subset\Omega$ there is a constant $C\ge0$ with
\[
  \bigl\|D^ku(x)-D^ku(y)\bigr\|\le C\,|x-y|^\alpha
  \qquad\text{for all }x,y\in F .
\]
\end{enumerate}
Here $D^ku(x)$ denotes the $k$-th total derivative of $u$ at $x$, a symmetric
$k$-linear map on $\mathbb{R}^2$, and $\|\cdot\|$ its operator norm.

This is the space $C^{k,\alpha}_{\mathrm{loc}}(\Omega)$ of Section~1 of the
paper, which is described there as the space of complex-valued $u\in
C^k(\Omega)$ whose $k$-th order partial derivatives $D^\nu u$, $|\nu|=k$, are
locally H\"older continuous with exponent $\alpha$; condition (b) is the
finiteness, for each compact $F$, of the second seminorm in the displayed
norm $\|u\|_{C^{k,\alpha}(F)}$ of Section~1.  Using the operator norm of the
$k$-th total derivative in place of the sum
$\sum_{|\nu|=k}|D^\nu u(x)-D^\nu u(y)|$ defines the same space: the two
quantities bound one another up to factors depending only on $k$, because the
partial derivatives $D^\nu u(x)$ with $|\nu|=k$ are precisely the values of
$D^ku(x)$ on the $k$-fold tuples drawn from the basis $\{1,i\}$ of
$\mathbb{R}^2$, and condition (b) is unchanged when $C$ is multiplied by such
a factor.
\end{definition}
